## 深度补充：TD学习算法详解

### TD(λ)的资格迹机制

资格迹是TD(λ)的核心，它记录了每个状态-动作对在过去被访问的"痕迹"。痕迹的衰减由λ参数控制：

数学定义：
$$ E_t(s,a) = \gamma \lambda E_{t-1}(s,a) + \mathbf{1}(S_t=s, A_t=a) $$

TD(λ)更新公式：
$$ Q(s,a) \leftarrow Q(s,a) + \alpha \delta_t E_t(s,a) $$
其中 $\delta_t = r + \gamma Q(s',a') - Q(s,a)$ 是TD误差。

λ的影响：
- $\lambda=0$：退化为TD(0)，高偏差低方差
- $\lambda=1$：接近蒙特卡洛，低偏差高方差  
- $\lambda \in (0,1)$：两者之间的权衡

### n-step TD的控制变种

| 算法 | 更新公式 | 特点 |
|------|----------|------|
| TD(0) | $V(s_t) \leftarrow V(s_t) + \alpha [r_{t+1} + \gamma V(s_{t+1}) - V(s_t)]$ | 单步，高偏差低方差 |
| n-step TD | $V(s_t) \leftarrow V(s_t) + \alpha [G_t^{(n)} - V(s_t)]$ | n步回报，平衡偏差方差 |
| TD(λ) | $V(s_t) \leftarrow V(s_t) + \alpha \delta_t E_t(s)$ | 资格迹，理论最优λ exists |

其中n步回报：
$$ G_t^{(n)} = r_{t+1} + \gamma r_{t+2} + \cdots + \gamma^{n-1} r_{t+n} + \gamma^n V(s_{t+n}) $$

### 代码示例：TD(λ)实现

```python
import numpy as np

class TDLambdaAgent:
    """TD(λ)实现（状态价值版本）"""
    def __init__(self, n_states, lamda=0.9, gamma=0.99, lr=0.01):
        self.V = np.zeros(n_states)
        self.lamda = lamda
        self.gamma = gamma
        self.lr = lr
    
    def update(self, trajectory, rewards):
        """
        trajectory: [s0, s1, ..., s_T]
        rewards: [r1, r2, ..., r_T]
        """
        T = len(trajectory)
        E = np.zeros_like(self.V)  # 资格迹
        
        for t in range(T):
            s_t = trajectory[t]
            
            # 计算TD误差
            if t < T-1:
                td_target = rewards[t] + self.gamma * self.V[trajectory[t+1]]
            else:
                td_target = rewards[t]  # 终止状态
            
            td_error = td_target - self.V[s_t]
            
            # 更新资格迹
            E *= self.gamma * self.lamda
            E[s_t] += 1.0
            
            # 更新所有状态（根据资格迹）
            self.V += self.lr * td_error * E
```

### 应用场景扩展

**1. 机器人导航（实时）**
- 机器人需要实时更新位置价值，不能等待episode结束
- TD学习可以每步更新，边走边学
- 实际案例：扫地机器人路径学习、无人机避障

**2. 股票交易（高频）**
- 每笔交易后需要立即更新策略，不能等收盘
- TD学习使用bootstrap，可以快速适应市场变化
- 实际案例：高频交易算法、加密货币套利

**3. 游戏AI（实时策略）**
- RTS游戏中需要实时决策，不能等一局结束
- TD学习可以边玩边学，持续优化策略
- 实际案例：星际争霸AI、文明AI

### 更多练习题

**练习3：数学推导**
问题：证明TD(0)在表格型情况下的更新等价于随机梯度下降在某个损失函数上。

**答案**：
考虑损失函数 $L(V) = \mathbb{E}[(R_{t+1} + \gamma V(S_{t+1}) - V(S_t))^2$
梯度：$\nabla_V L = -2(R_{t+1} + \gamma V(S_{t+1}) - V(S_t)) \cdot \nabla_V V(S_t)$
由于 $V(S_t)$ 对 $V(S_t)$ 的导数是1，对 $V(S_{t+1})$ 的导数是0（半梯度，不考虑bootstrap的影响），
因此 $\nabla_V L = -2\delta_t \cdot \mathbf{1}(S_t)$ （$\mathbf{1}(S_t)$是one-hot向量）
SGD更新：$V \leftarrow V - \frac{1}{2} \alpha \nabla_V L = V + \alpha \delta_t \cdot \mathbf{1}(S_t)$
这正是TD(0)的更新！

**练习4：λ参数选择**
问题：在某个具体任务中，如何选择λ？

**答案**：
1. 先用多个λ值（如0, 0.3, 0.5, 0.7, 0.9, 1.0）进行交叉验证
2. 绘制性能vs λ的曲线，选择最优λ
3. 通常：episode短选大λ（接近MC），episode长选小λ（接近TD(0)）
